%!TEX program = xelatex
\documentclass[11pt,a4paper]{ctexart}

% ==========================================
% 宏包引入与排版配置 (参考 RoboCup HSL Rules.tex 体系)
% ==========================================
\usepackage{geometry}
\geometry{left=2.5cm, right=2.5cm, top=2.8cm, bottom=2.8cm}

\usepackage{times}
\usepackage{fancyhdr}
\usepackage{url}
\usepackage{xcolor}
\usepackage{graphicx}
\usepackage{float}
\usepackage{amsmath,amssymb}
\usepackage{enumitem}
\usepackage{booktabs}
\usepackage{tabularx}
\usepackage{longtable}
\usepackage{array}
\usepackage{tcolorbox}
\tcbuselibrary{skins,breakable}

\usepackage[
    colorlinks=true,
    linkcolor=blue!70!black,
    urlcolor=blue!70!black,
    citecolor=green!50!black,
    bookmarksnumbered=true,
    pdfstartview=FitH
]{hyperref}

% 提示框宏定义 (仿照 Rules.tex hintbox 风格)
\newtcolorbox{noticebox}[1][]{
    colback=blue!5!white,
    colframe=blue!60!black,
    fonttitle=\bfseries,
    title=#1,
    arc=3mm,
    boxrule=1pt,
    left=4mm, right=4mm, top=3mm, bottom=3mm,
    breakable
}

\newtcolorbox{dangerbox}[1][]{
    colback=red!5!white,
    colframe=red!70!black,
    fonttitle=\bfseries,
    title=#1,
    arc=3mm,
    boxrule=1pt,
    left=4mm, right=4mm, top=3mm, bottom=3mm,
    breakable
}

\newtcolorbox{quotebox}{
    colback=gray!8!white,
    colframe=gray!40!white,
    arc=2mm,
    boxrule=0.5pt,
    left=4mm, right=4mm, top=2.5mm, bottom=2.5mm,
    breakable
}

% ==========================================
% 页眉与页脚定制 (每页标明 Robot工作坊 与 BigHeroX)
% ==========================================
\setlength{\headheight}{14pt}
\pagestyle{fancy}
\fancyhf{}
\lhead{\textbf{\textcolor{blue!80!black}{Robot工作坊}}}
\rhead{\textbf{\textcolor{red!80!black}{BigHeroX 战队}}}
\chead{\textcolor{gray}{\small RoboCup HSL 新人入门手册}}
\cfoot{\thepage}
\renewcommand{\headrulewidth}{0.8pt}
\renewcommand{\footrulewidth}{0.4pt}

% ==========================================
% 文档信息元数据
% ==========================================
\hypersetup{
    pdftitle={RoboCup HSL 新人入门手册 - Robot工作坊 BigHeroX},
    pdfauthor={hongli (2868014836@qq.com)},
    pdfsubject={RoboCup Humanoid Soccer League Onboarding Handbook},
    pdfkeywords={RoboCup, HSL, BigHeroX, RobotWorkShop, ROS 2, T4, Zvalley}
}


% ==========================================
% 内容标记：事实来源三种，写作规范见 docs/style-guide.md
% ==========================================
\newcommand{\official}{\textcolor{blue!70!black}{\footnotesize\textsf{[官方]}}}
\newcommand{\practice}{\textcolor{green!45!black}{\footnotesize\textsf{[本队]}}}
\newcommand{\advice}{\textcolor{orange!85!black}{\footnotesize\textsf{[建议]}}}
\newcommand{\todo}[1]{\textcolor{red!70!black}{\footnotesize\textsf{#1}}}
\newcommand{\pitfall}[1]{\begin{noticebox}[#1]}
\setcounter{tocdepth}{2}

\begin{document}

% ==========================================
% 封面排版 (Cover Page)
% ==========================================
\begin{titlepage}
    \centering
    \vspace*{0.5cm}
    
    % 双 Logo 居中展示
    \begin{figure}[h]
        \centering
        \includegraphics[height=3.0cm]{../public-resource/ima/RobotWorkShop.jpg}
        \hspace{1.5cm}
        \includegraphics[height=3.0cm]{../public-resource/ima/BigHeroX.jpg}
    \end{figure}
    
    \vspace{0.8cm}
    {\Huge \textbf{\textcolor{blue!90!black}{RoboCup 人形机器人足球联赛 (HSL)}}}\\[0.4cm]
    {\LARGE \textbf{新成员系统入门手册与技术工程图谱}}\\[0.5cm]
    {\large \textsf{A Comprehensive Onboarding Handbook \& Engineering Roadmap for RoboCup HSL}}
    
    \vspace{1.2cm}
    \rule{0.85\textwidth}{1.2pt}
    \vspace{1.0cm}
    
    \begin{minipage}{0.75\textwidth}
        \large
        \textbf{编著者（Author）：} hongli\\[0.3cm]
        \textbf{贡献者（Contributor）：} hanxu chen\\[0.3cm]
        \textbf{官方邮箱（Email）：} \href{mailto:2868014836@qq.com}{2868014836@qq.com}\\[0.3cm]
        \textbf{所属机构与战队：} 湖南大学 Robot 工作坊 · BigHeroX 战队\\[0.3cm]
        \textbf{开源仓库：} \url{https://github.com/RobotWorkShop-BigHeroX/RoboCupBook}\\[0.3cm]
        \textbf{手册版本：} \textsf{v1.2.0}\\[0.3cm]
        \textbf{构建日期：} \today
    \end{minipage}
    
    \vfill
    
    \begin{quotebox}
        \centering
        \small
        \textit{“面向对 RoboCup 及人形机器人足球联赛完全不了解的新成员。\\
        重点在建立清晰的技术知识框架、工程严谨思维与实验室安全底线。”}
    \end{quotebox}
    
    \vspace{0.8cm}
\end{titlepage}

% ==========================================
% 目录
% ==========================================
\thispagestyle{empty}
\tableofcontents
\newpage
\setcounter{page}{1}

% ==========================================
% 第 1 章：RoboCup 与 HSL
% ==========================================
\section{RoboCup 与 HSL}

RoboCup 于 1997 年举办首届赛事。其长期目标之一，是推动全自主机器人足球发展到能够与人类顶级球队竞技的水平。人形机器人足球是其中一条技术路线，要求机器人在没有人工遥控的情况下完成感知、决策和运动控制。\official

在这项比赛中，「全自主」意味着机器人需要独立完成比赛决策；「人形」则要求它在双足支撑和运动过程中维持平衡。感知、决策与运动控制不仅要分别工作，还必须在有限的计算资源和时间内协同运行。

\subsection{机器人足球考验什么}

机器人足球同时考验视觉、定位、策略和运动控制。机器人需要从相机图像中识别球与场地标记，估计自身位置，再根据比赛状态选择动作，最后通过关节控制完成移动或射门。这个过程会持续重复，任何一环的延迟或错误都可能传递到后续模块。例如，球的位置估计落后于真实位置时，控制器即使准确执行了目标动作，机器人也可能朝球已经离开的位置移动。

团队通常按这条链路划分工作：感知负责估计机器人、球和场地的位置；策略根据这些信息与比赛状态选择行为；运动控制将目标行为转换成可执行的关节动作。第 4 章按相关技术展开，读者可以先阅读与当前分工直接相关的部分。

比赛规则对自主运行和外部通信有明确限制。软件设计时应确认哪些信息可以从场外接收，并保证比赛策略能够在机载计算机上独立运行。具体限制以当届规则为准。

\subsection{相关组别与技术积累}

RoboCup 从来不止一个联赛。早些年同时铺着五六个方向，各自解决一块问题；最近几年，资源明显向人形足球收拢。

\begin{table}[H]
\centering
\small
\begin{tabularx}{\textwidth}{p{2.6cm} p{3.4cm} X}
\toprule
\textbf{组别} & \textbf{现在的状态} & \textbf{留给 HSL 的东西} \\
\midrule
SPL 标准平台组 & 并入人形体系 & 定位、队伍通信、多机战术的开源代码。B-Human 是这一支留下的完整遗产。 \\
\addlinespace
MSL 中型组 & 停办 & 全向底盘的运动规划和高速击球。双足用不上，但工程方法可借。 \\
\addlinespace
仿真组 & 压缩规模 & 便宜的训练与回归测试环境。它能替代真机跑通流程，替代不了真机调参数。 \\
\addlinespace
@Home 服务组 & 独立发展 & 家庭服务场景，与足球无关。本书不涉及。 \\
\bottomrule
\end{tabularx}
\caption{本书写作时的组别格局。调整每年都有，以当届官网公告为准。}
\end{table}

对开发者而言，组别变化带来的直接影响，是可参考的代码来自不同平台和赛制。人形机器人硬件、接触条件和场地光照都与轮式平台或纯仿真环境不同，因此复用旧方案时，需要先区分通用的软件设计与依赖具体硬件的实现。\advice

阅读早期开源项目时，可以优先研究定位、通信和状态机等系统设计，再判断哪些实现适用于当前平台。步态和平衡控制与硬件结构、执行器特性密切相关，通常不能直接移植。

\subsection{一场比赛的全过程}

参赛流程不仅包括场上的对抗，也包括检录、比赛状态切换和赛后处理。软件需要为每个阶段定义明确行为，尤其要保证机器人在非比赛状态下不会执行不应发生的动作。

\begin{table}[H]
\centering
\small
\begin{tabularx}{\textwidth}{p{2.4cm} X X}
\toprule
\textbf{阶段} & \textbf{场上发生什么} & \textbf{你的代码该做什么} \\
\midrule
赛前检录 & 裁判量尺寸、称重、查传感器，确认合规 & 提前核对尺寸与重量，并为必要的调整留出时间 \\
\addlinespace
INITIAL & 机器人在场外待命，等待入场 & 只收包，不动。所有电机保持安全状态 \\
\addlinespace
READY & 机器人自行走到开球站位 & 按规则到达指定位置，并确认定位状态满足启动条件 \\
\addlinespace
SET & 全部就位，静止 & 冻结底盘，准备开球动作 \\
\addlinespace
PLAYING & 自由对抗 & 感知、决策、运控全速跑；只有这个状态在计时 \\
\addlinespace
半场 & 交换场地，短暂休息 & 坐标整体镜像；重新初始化定位 \\
\addlinespace
FINISHED & 比赛结束 & 按规定进入安全状态，等待工作人员处理 \\
\bottomrule
\end{tabularx}
\caption{一场比赛的阶段与代码的对应关系。状态名和裁判盒一致，见 1.6 节。}
\end{table}

在团队赛制下，机器人需要分配站位、抢球和补位等任务，再根据队友状态进行协调。具体人数与换人规则可能随赛季调整，软件应从配置或裁判状态中读取相关信息，而不是把人数写死在逻辑中。\official

\subsection{场地条件对算法的影响}

\begin{table}[H]
\centering
\small
\begin{tabularx}{\textwidth}{p{2.2cm} p{4.2cm} X}
\toprule
\textbf{场地元素} & \textbf{它是什么} & \textbf{对算法意味着什么} \\
\midrule
草皮 & 人工草，表面可能反光 & 检查白平衡和曝光；如使用颜色阈值，应验证其对场地变化的敏感程度 \\
\addlinespace
白色标线 & 边界、中线、中圈、禁区 & 定位的主要特征来源；线宽固定，可用来估距离 \\
\addlinespace
中圈与罚球点 & 圆形和点状标记 & 圆形对称，容易产生镜像歧义，粒子滤波会收敛到错误一侧 \\
\addlinespace
球门 & 两根立柱加横梁 & 射门目标；立柱细，远了就丢，需要专门的检测策略 \\
\addlinespace
球 & 白底黑斑 & 滚动时外观变化大，检测要抗运动模糊 \\
\addlinespace
光照 & 场馆顶灯，也可能有自然光进入 & 光照可能随时间和场地位置变化，视觉参数应便于检查和调整 \\
\bottomrule
\end{tabularx}
\caption{场地元素、感知特征及其对算法的影响。具体尺寸以当届规则为准。}
\end{table}

现场出现感知或定位问题时，可以先根据表中的对应关系缩小排查范围，再用图像、日志和实测结果确认原因。

\subsection{规则里和写代码有关的六条}

以下列出几类会直接影响软件设计的规则要求。这份概述不能替代当届规则原文，赛前应逐项核对。

\begin{enumerate}[leftmargin=*]
    \item \textbf{全自主。} 比赛期间的自主运行、外部通信和人工干预均受规则约束。应确认程序在允许的机载环境中完成决策，并按队内调试规范保存日志。\official
    \item \textbf{尺寸与重量有上限。} 改机械之前先量一遍。检录不合格，整队白来。
    \item \textbf{传感器有限制。} 激光雷达明确禁用；深度相机在强光下噪点多，检录时属于灰色地带。定位只能靠可见光加里程计，第 4.7 节讲怎么做。
    \item \textbf{换人有代价。} 死球时才能申请，回场要等罚时结束。决策层要能处理「少一个人」的局面，也要能从零重新定位。
    \item \textbf{危险动作会被判罚。} 冲撞、推挤以及守门员持球时间等行为可能受到规则限制。策略层应将相关约束纳入行为选择，并在必要时允许机器人停止接近或主动退让。
    \item \textbf{半场交换场地。} 上下半场换边。代码里进攻方向不能写死，要能整体镜像一次。
\end{enumerate}

传感器限制会影响感知方案的选择；半场交换场地则要求定位、进攻方向和行为策略使用一致的坐标变换约定。

\subsection{裁判盒：机器人听谁的话}

裁判通过 GameController（比赛控制程序，队内简称裁判盒或 GC）广播当前比赛状态。机器人接收并解析这些信息，再据此决定是否允许移动以及何时执行比赛策略。

\begin{itemize}[leftmargin=*]
    \item \textbf{它怎么工作：} 主裁判在笔记本上运行 GC，选择比赛状态；GC 通过局域网把状态包广播到 3838 端口；场上所有机器人同时收到，各自解析。\official
    \item \textbf{它不管什么：} 不知道机器人在哪、球在哪。它只负责计时、计分、发状态、记罚时。你的定位、视觉、决策，它一概不参与。
    \item \textbf{它为什么重要：} 如果机器人收不到或错误解析状态包，就可能在不正确的比赛状态下继续运行。因此，集成测试应覆盖丢包、编号不一致和状态切换等情况。
\end{itemize}

\begin{table}[H]
\centering
\small
\begin{tabularx}{\textwidth}{p{2.4cm} X}
\toprule
\textbf{数据包字段} & \textbf{用途} \\
\midrule
队伍编号、球员编号 & 判断这个包属于哪台机器人；编号错了就全盘收不到 \\
\addlinespace
比赛状态 & INITIAL / READY / SET / PLAYING / FINISHED，以及各种罚时子状态 \\
\addlinespace
剩余时间 & 决定要不要抢最后一脚；也用来判断半场是否结束 \\
\addlinespace
罚时信息 & 谁被罚下、还剩多久，决定场上阵型怎么补 \\
\bottomrule
\end{tabularx}
\caption{裁判盒数据包里和决策相关的字段（字段名以官方仓库的定义为准）。}
\end{table}

\noindent\textbf{首次联调。} 建议按以下顺序逐项确认：

\begin{enumerate}[leftmargin=*]
    \item 下载官方 GameController，选对组别和队伍数量。
    \item 设好 Team Number 和 Player ID，和代码配置里逐位核对。
    \item 电脑和机器人放同一网段，关掉防火墙和虚拟网卡。
    \item 抓包确认：能收到广播，字段和时间都对得上。
    \item 走一遍完整的 INITIAL \(\rightarrow\) READY \(\rightarrow\) SET \(\rightarrow\) PLAYING，看机器人状态跳转。
\end{enumerate}

\noindent\textbf{三个常见坑：}

\begin{itemize}[leftmargin=*]
    \item \textbf{队号或球员号不一致。} 这类配置错误可能导致机器人忽略收到的数据包。联调时应先核对官方配置与程序中的编号。
    \item \textbf{防火墙或多网卡。} 网络配置可能阻止广播包到达程序，表现为间歇性丢包或只有部分机器人能接收。排查时应确认实际使用的网卡、网段和防火墙规则。
    \item \textbf{半场交换后方向错误。} 如果坐标变换不一致，机器人可能朝错误的球门移动。应在统一的坐标约定下处理换边，不要在多个模块中分别修改方向常量。
\end{itemize}

\subsection{五个常见误解}

\begin{description}[leftmargin=0pt,style=nextline]
    \item[「机器人是遥控的吧？」] 比赛要求机器人自主运行，人工干预受到规则限制。场边手持的控制器可能用于安全停机或进入特定保护状态，具体功能应按赛事规则和设备说明确认。
    \item[「必须买官方指定的机器人？」] 规则只要求合规：尺寸、重量、传感器都在规则里写着。自己造一台能过检录，完全可以参赛。
    \item[「仿真跑通了，真机应该也行。」] 仿真通过只能说明程序在当前模型和参数下表现正常。地面摩擦、执行器延迟、结构形变和光照变化都可能造成仿真与真机之间的差异，因此仍需逐步进行实机验证。
    \item[「规则是裁判的事，跟写代码没关系。」] 规则可能影响传感器选型、比赛状态处理、换人策略和坐标变换，因此需要在软件设计阶段一并考虑。
    \item[「代码写得越复杂越强。」] 复杂度本身并不代表性能优势。模块边界清楚、故障可定位、行为可复现，往往更有利于现场调试和长期维护。
\end{description}

\newpage
% ==========================================
% 第 2 章：强队与开源资产
% ==========================================
\section{强队与开源资产}

开源代码和技术报告可以帮助新成员了解成熟系统的组织方式，也能提供可供验证的实现参考。不过，不同赛季、规则和硬件之间存在差异，阅读时需要判断哪些设计可以复用。

下表列出几支值得研究的队伍，以及各自可能提供的参考方向。队伍成绩、技术报告和仓库状态应以官方资料为准。\official

\subsection{八支队伍及其参考方向}

\begin{table}[H]
\centering
\footnotesize
\begin{tabularx}{\textwidth}{p{1.9cm} p{2.6cm} p{1.5cm} X}
\toprule
\textbf{队伍} & \textbf{机构} & \textbf{地区} & \textbf{值得看的部分} \\
\midrule
B-Human & 不来梅大学 / DFKI & 德国 & 工程完整度。从相机到关节的整条链路都有文档，适合当第一份通读材料。 \\
\addlinespace
Rhoban FC & 波尔多大学 / CNRS & 法国 & 步态与踢球动作的参数化。想把动作调顺，看它的曲线设计。 \\
\addlinespace
NimbRo & 波恩大学 & 德国 & 全身动力学与自研硬件。要看「机器人怎么造」，看这里。 \\
\addlinespace
CIT Brains & 千叶工业大学 & 日本 & 高刚度关节控制与快速射门。 \\
\addlinespace
Hephaestus & 清华大学 & 中国 & 位姿估计与模型轻量化部署，和国内硬件适配的经验。 \\
\addlinespace
Mountain Sea & 中国农业大学 & 中国 & 多传感器融合、草皮上的步态补偿。 \\
\addlinespace
RoMeLa & UCLA & 美国 & 开源硬件平台（DARwIn-OP、ARTEMIS）的设计思路。 \\
\addlinespace
HERoEHS & 汉阳大学 & 韩国 & 结构轻量化与整机安全性设计。 \\
\bottomrule
\end{tabularx}
\caption{八支队伍与各自的可读之处。完整名单与近年战绩见 hsl.robocup.org/qualified-teams。}
\end{table}

\subsection{B-Human：从它的文档入手}

B-Human 会在赛后公开代码和技术报告。项目规模较大，初次阅读时不宜从目录中随机挑文件；先了解系统的功能分层，再沿着一条数据流追踪实现，通常更容易建立整体认识。

\begin{table}[H]
\centering
\small
\begin{tabularx}{\textwidth}{p{3.2cm} X}
\toprule
\textbf{层次} & \textbf{负责什么，你该看什么} \\
\midrule
感知（Perception） & 从图像里提特征：球、线、球门柱、其他机器人。看它怎么把一帧图拆成结构化的观测。 \\
\addlinespace
建模（Modeling） & 把观测变成状态：球的位置和速度、自己的位姿、对手在哪。看它的滤波器怎么写。 \\
\addlinespace
行为（Behavior） & 决定做什么：走位、抢球、传球、防守。看它的层次状态机怎么组织。 \\
\addlinespace
运动（Motion） & 把动作变成关节轨迹：步态、踢球、起身。看参数怎么暴露出来。 \\
\addlinespace
硬件（Hardware） & 和电机、IMU 通信。通常不用改，但要看得懂接口。 \\
\bottomrule
\end{tabularx}
\caption{B-Human 的分层结构。具体目录名以当年代码仓库为准。}
\end{table}

\noindent\textbf{建议的阅读顺序：}

\begin{enumerate}[leftmargin=*]
    \item 先读当年的技术报告，把五层的关系画成一张图。
    \item 再读数据结构定义。看它用什么类型表示「一个球」「一个位姿」——这决定了后面所有模块的接口。
    \item 然后找入口：一帧图进来，第一个函数在哪；一个决策周期开始，从哪调用。
    \item 最后挑一个模块深读，比如只读「球检测」这一条线，从图像到观测输出走通。
\end{enumerate}

\noindent\textbf{值得关注的三个工程设计：}

\begin{itemize}[leftmargin=*]
    \item \textbf{黑板式的数据流。} 模块把结果写进共享结构，别的模块按需读，模块之间不直接互相调用。好处是加一个模块不用改一圈调用链。
    \item \textbf{统一的坐标约定。} 全场只有一个「机器人坐标系」的定义，所有模块都用它。你们队要是每个模块各自转一次坐标，迟早对不上。
    \item \textbf{参数外置。} 阈值、速度、角度全放配置文件，现场能改。改代码再编译，比赛现场来不及。
\end{itemize}

\subsection{其他队伍的参考方向}

\begin{itemize}[leftmargin=*]
    \item \textbf{Rhoban：} 其踢球动作采用参数化曲线组织。研究这类实现时，可以关注动作参数如何暴露给调试人员，以及参数变化如何影响动作轨迹。
    \item \textbf{NimbRo：} 自研硬件与整机系统值得关注，尤其是结构设计、散热、布线和全身动力学之间的关系。
    \item \textbf{Hephaestus 与 Mountain Sea：} 可重点查阅它们公开的定位、感知和步态相关资料，并结合本队使用的平台与场地条件判断其适用性。\advice
\end{itemize}

\subsection{如何阅读陌生代码仓库}

\begin{enumerate}[leftmargin=*]
    \item \textbf{先看文档和架构图。} 通过 README、技术报告和模块图了解项目的目标、依赖与整体结构，再进入具体代码。
    \item \textbf{定位数据入口。} 确认图像、裁判盒状态和日志分别从哪里进入系统。入口明确后，才能沿数据流找到负责处理它们的模块。
    \item \textbf{追踪一条完整功能链。} 例如从图像中的球检测开始，继续查看定位、目标选择、接近球和踢球动作分别由哪些模块完成。
    \item \textbf{做一次小范围验证。} 修改一个参数，在相同条件下重新运行并比较结果，以确认自己对相关模块的理解。
\end{enumerate}

\noindent 代码阅读需要通过运行结果来验证。即使只调整一个参数，也应记录修改前后的现象，避免把对代码的猜测当成已经确认的结论。

\subsection{TDP 怎么读}

TDP（Team Description Paper，队伍技术说明文档）通常总结当年系统设计、实现方案和实验结果，可从官方队伍名单页获取。它面向评审，不一定包含所有实现细节，适合用于了解整体方案，再与代码、视频和历年版本对照。

阅读时可以优先核对以下三件事：

\begin{enumerate}[leftmargin=*]
    \item \textbf{他们用什么传感器、什么频率？} 这决定了方案的算力门槛，也决定了你能不能照搬。
    \item \textbf{他们的瓶颈写在哪？} 注意报告是否说明了系统的局限，例如远距离检测不稳定、定位易受对称结构影响等。明确的局限有助于理解方案的适用范围。
    \item \textbf{哪些内容是当年新增的？} 与上一年的报告对照，可以看出团队在哪些模块投入了开发资源，以及改动是否带来了可量化的结果。
\end{enumerate}

\noindent 技术报告呈现的是队伍选择记录和展示的内容。对于关键设计与性能结论，最好再查看实现、测试条件和对应版本。\advice

\newpage
% ==========================================
% 第 3 章：机器人平台
% ==========================================
\section{机器人平台：T4 与 Zvalley 生态}

本章涉及的尺寸、重量、关节限位和电气参数应来自厂家手册或实测记录。暂时缺少的数据应明确标记为待补，不应根据其他平台的参数推测。

\subsection{T4 平台参数与初步检查}

\begin{table}[H]
\centering
\small
\begin{tabularx}{\textwidth}{p{2.6cm} p{3.6cm} X}
\toprule
\textbf{参数} & \textbf{去哪查} & \textbf{为什么你必须知道} \\
\midrule
身高、体重 & 厂家手册（私有子模块） & 检录的硬指标；重心高低决定步态参数 \\
\addlinespace
自由度与限位 & SDK 文档、实物 & 写逆运动学之前必须知道每个关节能转到哪 \\
\addlinespace
关节电机与减速比 & 厂家手册 & 决定扭矩上限；踢球和起身容易超限 \\
\addlinespace
计算平台 & 开机查 \texttt{uname -a}、\texttt{nvidia-smi} & 决定你能跑多大的模型、几个进程 \\
\addlinespace
传感器清单 & 手册 + 实物 & 决定定位方案（见 4.7），也决定检录能不能过 \\
\addlinespace
电池与续航 & 手册 + 实测 & 决定一场比赛要换几次电池、什么时候换 \\
\bottomrule
\end{tabularx}
\caption{T4 参数与查阅途径。具体数值以厂家手册为准，本书不复述。}
\end{table}

\noindent\textbf{核对手册与实物。} 首次接手平台时，建议完成以下测量，并将结果保存到队内平台档案中：

\begin{enumerate}[leftmargin=*]
    \item 把机器人摆成标准零位，用卷尺量身高，用体重秤称整机。
    \item 逐个关节手动转到限位，记录正负方向的最大角度。
    \item 开机跑 \texttt{nvidia-smi} 和 \texttt{top}，记下空闲时的 CPU、内存、显存占用。
    \item 满电跑一次标准步态，记下电压下降曲线和可用的连续时间。
    \item 把上面四项写进一页纸，贴在工位。下一届直接接着用。
\end{enumerate}

\subsection{关节与自由度}

双足人形机器人的每条腿通常有 6 个自由度：髋部 3 个（前后摆、侧摆和旋转）、膝部 1 个、踝部 2 个（前后和侧向）。这些关节共同决定腿部的支撑、摆动和落脚能力，具体配置仍应以实际平台为准。

\begin{table}[H]
\centering
\small
\begin{tabularx}{\textwidth}{p{2.4cm} p{3.4cm} X}
\toprule
\textbf{部位} & \textbf{典型自由度} & \textbf{坏了会怎样} \\
\midrule
髋（前后） & 1 & 迈不开步，步幅直接减半 \\
髋（侧向） & 1 & 站不稳，重心移不过去 \\
髋（旋转） & 1 & 转身别劲，踢球方向偏 \\
膝 & 1 & 蹲不下去，起身动作做不出来 \\
踝（前后） & 1 & 缓冲失效，容易前扑 \\
踝（侧向） & 1 & 侧向抗扰差，被撞就倒 \\
\bottomrule
\end{tabularx}
\caption{单腿各关节自由度的典型分工。具体数量与方向以实物为准。}
\end{table}

\begin{noticebox}[注意]
关节零位由机械标定确定，编码器读数需要与标定结果保持一致。更换电机或编码器、发生较严重的跌倒，或发现关节位置异常后，应按厂家规定检查并重新标定。具体流程以平台手册为准。
\end{noticebox}

\subsection{IMU 与姿态估计}

IMU（惯性测量单元）通常提供角速度、加速度以及由内部或外部算法估计的姿态信息：

\begin{itemize}[leftmargin=*]
    \item \textbf{三轴角速度}：转身、晃动有多快。
    \item \textbf{三轴加速度}：重力方向在哪，有没有被撞。
    \item \textbf{姿态角}：芯片内部解算好的横滚、俯仰、偏航。
\end{itemize}

在控制系统中，IMU 常用于姿态反馈和跌倒检测。姿态估计结果可以用于调整踝关节与髋关节的控制目标；当姿态超过设定范围时，保护逻辑可触发相应的安全动作。

\begin{table}[H]
\centering
\small
\begin{tabularx}{\textwidth}{p{2.6cm} X}
\toprule
\textbf{坑} & \textbf{现象与处理} \\
\midrule
零偏 & 静止时角速度不为零，积分几秒就漂出几度。上电后先静止采样几秒求均值，减掉再积分。 \\
\addlinespace
温漂 & 传感器温度变化时零偏可能发生变化。可在规定的预热条件下检查，必要时建立温度补偿模型。 \\
\addlinespace
振动 & 走路时加速度全是高频毛刺，姿态跟着抖。加低通滤波，截止频率从 30 Hz 往下试。 \\
\addlinespace
坐标系 & 芯片坐标系和机体坐标系不一定重合。装的时候记下朝向，代码里做一次固定旋转。 \\
\bottomrule
\end{tabularx}
\caption{IMU 测量中常见的误差来源与处理方法。}
\end{table}

\subsection{足式里程计与误差来源}

轮式机器人靠轮子转速算位移。人形没有轮子，只能靠腿：判断哪只脚在地上（支撑相），用正运动学把关节角度换算成脚的位置，再反推身体相对移动了多少。

\begin{table}[H]
\centering
\small
\begin{tabularx}{\textwidth}{p{2.8cm} X}
\toprule
\textbf{误差来源} & \textbf{怎么表现，怎么缓解} \\
\midrule
打滑 & 实际位移小于关节运动所估计的位移。可结合视觉标线或其他外部观测修正里程计误差。 \\
\addlinespace
足底柔性 & 脚掌受力形变，算出来的触地位置比实际偏。用触地传感器做校正。 \\
\addlinespace
关节零位误差 & 零位差一度，走十米就偏出几十厘米。定期标定。 \\
\addlinespace
时间不同步 & 图像和关节数据的采集时刻差几十毫秒，融合时对不上。统一时间戳。 \\
\bottomrule
\end{tabularx}
\caption{足式里程计的常见误差来源及缓解方法。}
\end{table}

因此，足式里程计适合提供连续的短时相对运动估计，但误差会随时间累积；视觉观测可用于约束长期漂移。两者的融合方法见第 4.7 节。

\subsection{机载软件结构与运行频率}

机载软件通常由多个进程或节点组成。不同模块的实时性要求不同，因此通常需要分别设置运行频率：

\begin{table}[H]
\centering
\small
\begin{tabularx}{\textwidth}{p{2.4cm} p{1.8cm} p{4.4cm} X}
\toprule
\textbf{模块} & \textbf{典型频率} & \textbf{输入} & \textbf{输出} \\
\midrule
相机驱动 & 30--60 Hz & 图像 & 原始帧 \\
视觉检测 & 15--30 Hz & 原始帧 & 球、球门、机器人的位置 \\
定位 & 10--30 Hz & 视觉观测、里程计 & 自己在场地的位姿 \\
决策 & 5--20 Hz & 定位、球、裁判状态 & 目标动作 \\
运动控制 & 100--500 Hz & 目标动作、IMU & 关节指令 \\
\bottomrule
\end{tabularx}
\caption{机载软件各模块的典型运行频率。实际值应通过平台测量与配置确认。}
\end{table}

\subsection{电气系统与通信接口}

电源、总线和急停回路需要与软件接口一并检查。接线错误可能损坏驱动器或造成安全风险，因此电气操作应遵循平台手册和实验室流程。

\begin{table}[H]
\centering
\small
\begin{tabularx}{\textwidth}{p{2.4cm} p{4.6cm} X}
\toprule
\textbf{部分} & \textbf{它做什么} & \textbf{碰之前注意什么} \\
\midrule
动力母线 & 给关节电机供电，电压高、电流大 & 带电插拔会打火，接反直接烧驱动 \\
\addlinespace
逻辑电源 & 给工控机、传感器供电 & 和动力电分开走，避免电机启停把板子拉黑 \\
\addlinespace
CAN 总线 & 工控机和各关节通信 & 终端电阻要对，线太长会丢帧 \\
\addlinespace
以太网 & 相机、工控机、调试机互联 & 和裁判盒同网段，别接错交换机 \\
\addlinespace
急停回路 & 硬件切断动力 & 必须独立于软件，软件停止响应时它仍然有效 \\
\bottomrule
\end{tabularx}
\caption{机器人主要电气与通信接口及操作注意事项。}
\end{table}

\noindent\textbf{接线与上电检查：}

\begin{enumerate}[leftmargin=*]
    \item 插拔任何线之前，先断动力电。带电插拔 CAN 线和电机线，是烧驱动板的常见原因。
    \item 上电顺序固定：逻辑电先，自检通过，再上动力电。顺序反过来，工控机可能被浪涌损坏。
    \item 电池不过放。电压掉到保护阈值以下，电池就废了；长期存放要充到一半。
    \item 急停回路不接软件。它可以被软件触发，但断电的那条路径必须是硬的。
    \item 改线后先用万用表检查通断和短路情况，确认接线无误后再上电。
\end{enumerate}

\begin{noticebox}[经验]
排查驱动板故障时，应先核对最近一次接线和上电操作，再检查器件本身。故障原因需要通过测量和日志确认，不宜仅凭现象判断。\advice
\end{noticebox}
\subsection{队内平台与资料}

队里的私有子模块（\texttt{Zvalley-resource-old/}）放着三类东西：厂家硬件手册、历年比赛录像与裁判盒日志、内部培训讲义。克隆要带子模块参数：

\begin{verbatim}
git clone --recurse-submodules <仓库地址>
git submodule update --init --recursive
\end{verbatim}

\begin{noticebox}[权限]
私有子模块只对授权队员开放。原理图、固件协议和内部参数等资料应从该模块或平台负责人处获取；如果现有资料无法确认某项配置，应先向负责人核实。
\end{noticebox}

\subsection{Zvalley 生态：五个入口}

\begin{table}[H]
\centering
\footnotesize
\begin{tabularx}{\textwidth}{p{3.2cm} X}
\toprule
\textbf{入口} & \textbf{什么时候用} \\
\midrule
开放资料总站 & 找固件版本、接口说明、工具链更新 \\
\addlinespace
SDK 文档中心 & 查电机通信协议、总线读写、传感器驱动接口 \\
\addlinespace
SDK 仓库（GitHub） & 拿 C++ / Python 示例，先跑通一个关节再谈其他 \\
\addlinespace
Z01 仿真资产 & 要 URDF / MJCF / USD 模型，做仿真或在 Isaac 里验证 \\
\addlinespace
Zlab\_Gym 手册 & 想用强化学习训练步态，从这里起步（见 4.8） \\
\bottomrule
\end{tabularx}
\caption{Zvalley 官方资料与开发工具入口。网址见附录 F 或队内书签。}
\end{table}

\noindent\textbf{SDK 的初步验证流程：}

\begin{enumerate}[leftmargin=*]
    \item 克隆 SDK 仓库，按 README 装依赖。
    \item 跑官方示例，先只控制一个关节，确认能读能写。
    \item 打印所有关节的位置反馈，和实物角度对一遍。
    \item 写一个最小循环：读 IMU、读关节、发一个固定目标位置。
    \item 加入超时处理、角度限幅和急停回调等保护，并在受控条件下验证。保护逻辑未经确认前，不应驱动真机关节。
\end{enumerate}

\subsection{本队公开资料索引}

\begin{table}[H]
\centering
\small
\begin{tabularx}{\textwidth}{p{2.8cm} X p{2.6cm}}
\toprule
\textbf{材料} & \textbf{里面有什么} & \textbf{什么时候看} \\
\midrule
2025 TDP & 上一代方案：机械结构、视觉流程 & 想了解设计取舍 \\
2026 TDP & 2026 赛季方案：7 对 7 的分布式决策与步态 & 接手前的第一份材料 \\
展示视频 & 自主寻球、带球、射门、起身的完整流程 & 给新人建立直觉 \\
Robots.pdf & 电机选型、连杆尺寸、重心 & 改机械之前 \\
software.pdf & 软件架构、话题拓扑、状态转移 & 读代码之前 \\
\bottomrule
\end{tabularx}
\caption{本队公开资料。全部放在 \texttt{public-resource/} 下。}
\end{table}

\newpage
% ==========================================
% 第 4 章：学习路线
% ==========================================
\section{学习路线与开发基础}

\begin{quotebox}
\small
学习机器人开发时，不必一开始就深入所有底层细节。对于 ROS 2、Linux 与 Git 等工具，应当先理解它们要解决的工程问题，再依据所承担的实际开发任务逐步深入。\\
既不能完全依赖生成式工具输出未经推导的代码，也不应陷入漫无目的地逐行阅读依赖库的低效怪圈。最有效的起点，是先亲手跑通一条端到端的完整功能链路，再围绕遇到的具体故障定向深挖。\\
禁止完全许愿式编程，但也不必转向一行行了解代码，这就像从一个低效走向另一个低效。\\
\hfill —— andinm\_hong
\end{quotebox}

新成员可根据当前确定的研发方向选择阅读路径：

\begin{table}[H]
\centering
\small
\begin{tabularx}{\textwidth}{p{2.2cm} X p{4.2cm}}
\toprule
\textbf{研发方向} & \textbf{必读技术章节} & \textbf{可以先跳过的部分} \\
\midrule
感知组 & 4.4 Linux、4.5 YOLO 目标检测、4.7 粒子滤波定位，结合 3.3 IMU 与 3.4 里程计 & 4.8 深度强化学习数学推导 \\
\addlinespace
策略组 & 4.3 ROS 2 通信机制、4.6 行为树，结合 1.6 裁判盒状态机 & 4.7 粒子滤波内部积分数学推导 \\
\addlinespace
运控组 & 4.4 Linux、4.6 行为树动作接口，结合 3.2 关节自由度与 3.3 IMU 姿态解算 & 4.5 视觉检测网络训练与数据增强细节 \\
\bottomrule
\end{tabularx}
\caption{按研发方向规划的学习路径。公共基础为 Git 协作、Linux 命令行与 ROS 2 通信模型。}
\end{table}

\subsection{可以暂缓深入的内容}

之所以建议新人在入队初期暂缓钻研以下方向，是因为在尚未建立整机数据流全貌之前，过早陷入孤立的底层细节极易引发严重的挫败感，迫使学习进度停滞在细枝末节上：

\begin{itemize}[leftmargin=*]
    \item \textbf{DDS 底层报文协议细节。} 初期只需掌握节点自动发现、话题发布订阅与服务调用的基本使用，只有在遭遇局域网内广播风暴或跨网段断连等网络疑难杂症时，才需要抓包分析 RTPS 报文头。
    \item \textbf{Executor 调度器的内部源码。} 先掌握单线程与多线程执行器的调度差异，明确耗时计算会阻塞同一执行器内的其他回调即可，无需通读其源码调度循环。
    \item \textbf{C++ 高阶模板元编程。} 现有工程中大量运用了模板封装，阅读时能看懂接口调用形式即可，避免在语法技巧上耗费过多时间。
    \item \textbf{强化学习算法收敛性证明。} 重点理解训练流程、动作空间与观测空间的定义方式，数学证明不影响工程调试。
    \item \textbf{复杂文档排版技巧。} 集中精力于技术事实与工程逻辑的准确表达，排版工具按需查阅即可。
\end{itemize}

\subsection{Git：建立严格的工程协作习惯}

在多人协同开发大型机器人工程时，Git 不仅是版本备份工具，更是维护代码稳定性的第一道防线：

\begin{table}[H]
\centering
\small
\begin{tabularx}{\textwidth}{p{3.2cm} p{4.6cm} X}
\toprule
\textbf{任务场景} & \textbf{推荐命令} & \textbf{操作目的与注意事项} \\
\midrule
查看工作区变动 & \texttt{git status}、\texttt{git diff} & 提交前逐行比对修改，避免误将本地临时测试日志或中间产物混入暂存区 \\
\addlinespace
精细化暂存 & \texttt{git add -p} & 交互式逐块审查改动，杜绝将无关修改强行打包 \\
\addlinespace
原子化提交 & \texttt{git commit -m "..."} & 保持单次提交的语义原子性，便于后续故障二分回溯（\texttt{git bisect}） \\
\addlinespace
同步远程分支提交 & \texttt{git pull --rebase} & 维持清晰线性的提交历史，避免产生大量无意义的分支合并节点 \\
\addlinespace
开辟特性分支 & \texttt{git switch -c feat/...} & 隔离新功能开发，维持主分支始终处于可编译、可测试状态 \\
\addlinespace
丢弃工作区修改 & \texttt{git restore <file>} & 放弃未提交的临时改动；若需撤销已推送提交，必须使用 \texttt{git revert} 生成逆向提交 \\
\addlinespace
更新私有子模块 & \texttt{git submodule update ...} & 使底层依赖库与私有技术手册同步至指定版本校验码 \\
\bottomrule
\end{tabularx}
\caption{常用 Git 命令及其操作规程。}
\end{table}

团队协作中必须严格恪守三条纪律：
\begin{enumerate}[leftmargin=*]
    \item \textbf{严禁对公共受保护分支执行强行推送（\texttt{force push}）。} 强行覆盖远程分支将直接破坏其他协作者的本地提交基底，导致灾难性的代码丢失与繁琐的重组代价。
    \item \textbf{坚持单一职责原则。} 严禁在一个提交中同时混合“功能开发”与“全工程代码格式化”。这会导致代码审查完全无法聚焦，在遭遇缺陷需要回溯时，也会直接剥夺原子化撤销的可能性。
    \item \textbf{提交 PR 前必须在本地通过全套构建与静态检查。} CI 自动化工作流仅作为辅助防线，开发者必须对自身推送代码的本地编译与基本功能负责。
\end{enumerate}

\subsection{ROS 2 的常用机制与适用场景}

ROS 2 是机器人软件的神经系统。各个核心机制的设计直接对应着特定的工程解题需求：

\begin{table}[H]
\centering
\small
\begin{tabularx}{\textwidth}{p{2.6cm} p{5.0cm} X}
\toprule
\textbf{核心机制} & \textbf{解决的工程问题} & \textbf{常见调试与故障场景} \\
\midrule
DDS / RTPS & 去中心化节点自动发现与高吞吐网络传输 & 局域网内某台机器人重启后，排查节点是否能自动恢复拓扑连接 \\
\addlinespace
节点与通信模型 & 模块独立运行，通过话题、服务与动作交换数据 & 梳理系统架构与数据调用关系，划分模块职责边界 \\
\addlinespace
QoS 服务质量策略 & 规定消息可靠性、队列缓冲深度与保活要求 & 解决相机画面延迟过高、控制指令堆积或比赛状态偶发丢包 \\
\addlinespace
Executor 执行器 & 调度回调函数并发线程与优先级 & 解决控制周期抖动、高频通信丢帧与多回调死锁问题 \\
\addlinespace
Lifecycle 生命周期 & 提供明确的节点状态机（未配置/激活/失活） & 比赛上场前顺序自检与停机安全保护 \\
\addlinespace
Node Composition & 支持多节点在同一进程中运行，实现零拷贝通信 & 解决高分辨率高帧率图像在进程间拷贝占用过高 CPU 资源的问题 \\
\addlinespace
TF2 坐标变换系统 & 维护机器人机体、传感器与球场的空间树形转换 & 将图像中识别到的目标映射为球场绝对坐标 \\
\addlinespace
rosbag2 数据录制 & 完整记录现场传感器与话题通信时序数据 & 赛后离线回放现场数据复现偶发故障 \\
\bottomrule
\end{tabularx}
\caption{ROS 2 核心机制与实际调试场景对应关系。}
\end{table}

\noindent\textbf{QoS 配置策略的选择依据：}

\begin{itemize}[leftmargin=*]
    \item \textbf{视频图像流：} 相机传输 1080P/720P 原始帧时数据通量极高。必须采用 \texttt{BEST\_EFFORT}（尽力而为）传输模式并将队列深度设为 1。之所以这样设置，是因为在图像处理链路中，下一帧新图像永远比上一帧更具时效价值；如果误配置为 \texttt{RELIABLE}（可靠传输），一旦出现单帧网络丢包，通信中间件将不断发起重传重试，由此导致历史画面帧在缓冲区发生严重堆积，造成机器人看到的视野产生数秒的时延滞后。
    \item \textbf{裁判盒与状态机指令：} 必须采用 \texttt{RELIABLE} 模式。比赛状态切换指令低频但极为致命，漏收一次便会导致机器人无法按时开球或违规受罚，绝不允许丢失。
    \item \textbf{高频 IMU 姿态反馈：} 通常采用 \texttt{BEST\_EFFORT} 传输且队列设为 1，使运动控制回调函数获取到的数值始终保持当前最新。
\end{itemize}

\subsection{Linux 与终端开发基础}

Ubuntu Linux 是机器人机载主控的标准操作系统，熟练掌握命令行是执行真机联调的基础：

\begin{table}[H]
\centering
\small
\begin{tabularx}{\textwidth}{p{3.2cm} p{4.2cm} X}
\toprule
\textbf{命令分类} & \textbf{常用命令} & \textbf{实际工程应用场景} \\
\midrule
进程与资源监控 & \texttt{htop}、\texttt{ps aux} & 排查机载 CPU 占用是否过载、识别异常僵死进程 \\
\addlinespace
系统日志排查 & \texttt{journalctl -u ... -f} & 实时追踪开机自启系统服务的报错与崩溃追踪信息 \\
\addlinespace
磁盘空间排查 & \texttt{df -h}、\texttt{du -sh *} & 避免录制大量 rosbag2 数据导致工控机固态硬盘爆满死机 \\
\addlinespace
网络端口诊断 & \texttt{ip addr}、\texttt{netstat} & 确认调试网口 IP 配置、抓取裁判盒 3838 广播端口报文 \\
\addlinespace
软硬件权限管理 & \texttt{sudo usermod -aG ...} & 解决 USB 串口与 CAN 驱动板无权限读写的底层错误 \\
\bottomrule
\end{tabularx}
\caption{机载 Linux 系统维护核心命令。}
\end{table}

\subsection{YOLO 目标检测：从标注到推理}

视觉感知的核心目标是从图像像素中提取结构化特征。训练与部署 YOLO 必须遵循严谨的工程闭环：现场视频采集抽帧 $\rightarrow$ 标注足球与球门 $\rightarrow$ 数据增强与色温抖动 $\rightarrow$ 模型训练与 mAP 评估 $\rightarrow$ 导出 ONNX / TensorRT $\rightarrow$ 真机推理与时延评测。

\begin{enumerate}[leftmargin=*]
    \item \textbf{样本多样性压倒模型复杂度。} 之所以要在人工草皮现场采集多时间段样本，是因为训练数据集一旦缺乏强反光、强逆光或大面积阴影样本，即使选用超大规模的模型架构，在真实赛场上也会频繁发生漏检。
    \item \textbf{实时性是第一底线。} 视觉检测运行在机载 i7-1355U 上，算力资源十分紧张。检测耗时必须严格压缩在 30 毫秒以内（至少维持 30Hz 推理速度），否则将直接加剧控制回路的总时延。通常建议将 PyTorch 权重导出为 ONNX，并利用 TensorRT 或 OpenVINO 执行 FP16 量化加速。
\end{enumerate}

\subsection{行为树：任务规划与状态解耦}

战术策略层通常采用行为树（BehaviorTree）替代复杂的嵌套状态机：

\begin{table}[H]
\centering
\small
\begin{tabularx}{\textwidth}{p{2.4cm} p{3.2cm} X}
\toprule
\textbf{节点类型} & \textbf{典型节点} & \textbf{调度机制与工程作用} \\
\midrule
控制节点 & \texttt{Fallback} & 顺序尝试执行各子节点，只要有一个子节点返回 SUCCESS 便立刻返回，用于处理容错兜底 \\
\addlinespace
控制节点 & \texttt{Sequence（起身）} & 仅当所有子节点均返回 SUCCESS 时才返回成功，用于执行倒地起身固定序列 \\
\addlinespace
条件节点 & \texttt{IsFallen} & 仅读取 IMU 姿态判定机体是否倒地，只读查询，严禁引入改变外部状态的副作用 \\
\addlinespace
动作节点 & \texttt{GetUp} & 触发底层运动控制器执行起身轨迹 \\
\addlinespace
控制节点 & \texttt{Sequence（进攻）} & 顺序调度寻球、接近球体与射门动作 \\
\addlinespace
动作节点 & \texttt{FindBall / Kick} & 执行旋转寻球、调整身位逼近与单腿踢球动作 \\
\bottomrule
\end{tabularx}
\caption{行为树核心节点及其调度机制。}
\end{table}

编写行为树时必须遵守两项核心准则：
\begin{itemize}[leftmargin=*]
    \item \textbf{条件节点绝对严禁引入副作用。} 条件节点在行为树每个 tick 周期都会被高频轮询。如果开发者在检查“是否看到足球”的同时顺手修改了全局目标变量，当行为树发生逻辑回溯时，将引发不可预测的循环震荡。
    \item \textbf{严格匹配决策 tick 周期与底层控制频率。} 行为树通常以 10 至 20 Hz 的低频运行，而底层运动控制运行在 200 Hz。行为树下发的瞬态速度跳变必须经过平滑加速度插值器过滤，避免阶跃速度指令直接导致电机过流或机体失稳。
\end{itemize}

\subsection{粒子滤波与场地定位}

足式里程计能够推算短时位移，但打滑和累计误差会持续增大；前置相机能看到场地白线、中圈和球门，但因为足球场两头形状完全一致，容易认错半场。蒙特卡洛粒子滤波（MCL）在地图坐标系中维护一组候选位置点（粒子），用里程计推算它们怎么移动，再用看到的标线和球门给每个粒子打分，最后算出机器人在球场上的具体坐标。

\subsubsection{定位计算步骤}

\begin{enumerate}[leftmargin=*]
    \item \textbf{初始撒点：} 在球场范围内生成大量随机粒子，每个粒子记录一个候选坐标与朝向角度 $(x, y, \theta)$。
    \item \textbf{叠加运动推算：} 读取里程计位移增量，把机体走过的距离加到每个粒子上，并加入运动高斯噪声模拟打滑。
    \item \textbf{比对视觉标线：} 把相机检测到的白线交点、中圈弧线和球门立柱，与各粒子视角下应该看到的地图特征比对，重合度越高给的分数（权重）越高。
    \item \textbf{重采样筛选：} 淘汰得分过低的粒子，在得分高的粒子周围重新复制生成新粒子。
    \item \textbf{计算当前坐标：} 统计聚集在一起的高分粒子重心，输出机器人在球场坐标系中的当前位置与朝向。
\end{enumerate}

\begin{table}[H]
\centering
\small
\begin{tabularx}{\textwidth}{p{2.6cm} p{3.2cm} X}
\toprule
\textbf{调参项目} & \textbf{典型范围} & \textbf{调整影响与权衡取舍} \\
\midrule
粒子数量 & 500--2000 个 & 粒子过多拖慢机载工控机计算；粒子过少容易跟丢位置 \\
\addlinespace
运动噪声参数 & 根据实测打滑程度标定 & 设置过小会导致跟丢真实轨迹；设置过大会让粒子散得太开 \\
\addlinespace
观测噪声参数 & 根据相机标定精度设定 & 设置过小会过度轻信单次画面抖动；设置过大会忽略看到的有效标线 \\
\addlinespace
重采样阈值 & 有效粒子降到总数一半 & 重采样太频繁会丧失寻找备选位置的能力，太慢会导致坐标更新迟钝 \\
\bottomrule
\end{tabularx}
\caption{粒子滤波关键调优参数及其工程权衡。}
\end{table}

\noindent\textbf{常见定位问题与处理方法：}

\begin{itemize}[leftmargin=*]
    \item \textbf{半场认反。} 足球场关于中线完全对称，两头的白线和球门形状一模一样。如果机器人开场撒点不确定，或者被撞倒爬起后丢失了朝向，仅靠看到的白线就会把两个半场搞混，把自己的位置算到对方半场去。实际处理时需要结合机载指南针航向角锁定攻防方向，或者在开赛前根据裁判盒下发的队伍半场信息初始化朝向。
    \item \textbf{标线被遮挡。} 机器人身处开阔中场或视线被其他机器人挡住时，视野里找不到有效标线。此时依靠里程计做短时推算，暂停粒子重采样，避免在缺乏证据时把粒子胡乱收拢。
    \item \textbf{碰撞后重新定位。} 遭受猛烈冲撞导致机器人完全丢失当前位置时，重新在全场撒点，旋转机身寻找多个标线特征重新计算位置。
\end{itemize}

\subsection{深度强化学习的适用范围}

本节为进阶选读内容，梳理双足运动控制领域常见的基础概念：

\begin{table}[H]
\centering
\small
\begin{tabularx}{\textwidth}{p{3.2cm} p{4.6cm} X}
\toprule
\textbf{核心概念} & \textbf{数学物理内涵} & \textbf{在双足控制中的具体角色} \\
\midrule
贝尔曼方程 & 当前状态价值等于即时奖励加上经折减的未来期望价值 & 构成值函数迭代优化的理论基石 \\
\addlinespace
时序差分（TD） & 基于单步或多步转移样本更新当前价值函数 & 避免等待完整对抗回合结束，大幅提升样本利用效率 \\
\addlinespace
策略网络 & 以高频观测数据为输入，直接输出各关节目标动作 & 端到端映射状态到力矩，替代传统复杂的运动学解析解 \\
\addlinespace
Actor-Critic 架构 & Actor 网络生成动作，Critic 网络评估打分 & 降低策略梯度估计方差，促使训练平稳收敛 \\
\addlinespace
PPO 算法 & 严格限制每次策略更新的概率分布剪切步长 & 防止单次梯度更新幅度过大导致策略崩溃，在双足步态训练中应用最为广泛 \\
\bottomrule
\end{tabularx}
\caption{强化学习常见基础概念及其控制角色。}
\end{table}

\noindent\textbf{部署到真机前的物理制约与风险评估：}

\begin{itemize}[leftmargin=*]
    \item \textbf{真机跌落与机械磨损代价。} 双足机器人在真机上直接试错成本极高。每次高动态跌落都会对高减速比行星齿轮箱与轴承施加剧烈剪切载荷，因此策略训练与初步筛选必须在 Isaac Gym 或 MuJoCo 等高并发物理仿真引擎中完成。
    \item \textbf{仿真与现实差距（Sim2Real Gap）。} 仿真接触模型与实际人工草皮摩擦力、电机驱动器响应延迟存在差异。训练时必须引入域随机化（Domain Randomization）对质量分布、摩擦系数与通信延迟施加扰动。
    \item \textbf{黑盒策略可解释性。} 端到端神经网络难以直观分析特定动作选择的具体原因，现场调试与紧急修复难度较大。通常建议结合残差控制或基于经典控制框架进行策略融合。\advice
\end{itemize}

\newpage
% ==========================================
% 第 5 章：消息渠道
% ==========================================
\section{消息渠道}

规则修订、赛程调整和报名通知可能分散在多个官方渠道。建议将以下入口加入队内资料库，并明确由谁定期检查。

\subsection{优先关注的两个渠道}

\begin{enumerate}[leftmargin=*]
    \item \textbf{官方邮件列表。} 技术委员会和组委会可能通过邮件发布规则修订、资格通知和重要裁决。订阅后仍应留意当届官网的补充通知。\official
    \begin{itemize}
        \item 订阅地址：\url{https://lists.robocup.org/mm/lists/robocup-humanoid.lists.robocup.org/}
        \item 建议使用队内公共邮箱订阅，便于成员交接，避免个人离队后通知无人接收。
    \end{itemize}
    \item \textbf{官方 Discord。} 这里可用于了解赛事讨论、现场网络问题和规则解释等信息。回复速度取决于具体问题与在线人员情况。\official
    \begin{itemize}
        \item 加入链接：\url{https://discord.gg/67qGN8ycwz}
        \item 提问前先搜索历史消息，并附上相关日志、环境信息和复现步骤，以便其他队伍定位问题。
    \end{itemize}
\end{enumerate}

\subsection{官方信息入口}

\begin{table}[H]
\centering
\small
\begin{tabularx}{\textwidth}{p{4.6cm} X p{2.2cm}}
\toprule
\textbf{站点} & \textbf{上面有什么} & \textbf{多久看一次} \\
\midrule
hsl.robocup.org & 联赛主页：赛程、组别、技术挑战主题、入围名单 & 每月 \\
\addlinespace
202x.robocup.org & 当届主办方的现场信息：场馆、注册、网络、草皮 & 赛前每周 \\
\addlinespace
/information-for-new-teams & 新队指南：资格流程、TDP 规范、现场准备 & 报名前 \\
\addlinespace
/qualified-teams & 所有入围队伍的 TDP、视频、硬件配置 & 赛前必读 \\
\addlinespace
robotprogram.robocup.org & 官方机器人支持计划，样机与技术支持 & 想换平台时 \\
\bottomrule
\end{tabularx}
\caption{官方信息与资源入口。}
\end{table}

\subsection{代码与工具仓库}

\begin{table}[H]
\centering
\small
\begin{tabularx}{\textwidth}{p{3.8cm} X}
\toprule
\textbf{仓库} & \textbf{你会用到它的场景} \\
\midrule
HSL-Rules & 查规则原文；看某条规则的修订历史；确认某个限制是哪一年加上的 \\
\addlinespace
GameController & 本地联调裁判盒；对照源码确认数据包字段 \\
\addlinespace
RobotInspection & 赛前自检，提前知道哪些尺寸会被卡 \\
\bottomrule
\end{tabularx}
\caption{常用官方仓库。核对规则时，应确认文件版本与最近的修订记录。}
\end{table}

\noindent 这三个仓库位于 \url{https://github.com/RoboCup-HumanoidSoccerLeague}。追踪规则变化时，应确认 PDF 与仓库源文件对应的版本，并查看相关提交记录。

\subsection{一个赛季的节奏}

具体日期每年都会调整。下表概括一个赛季中常见的准备阶段，可用于安排队内任务：\advice

\begin{table}[H]
\centering
\small
\begin{tabularx}{\textwidth}{p{2.6cm} X}
\toprule
\textbf{阶段} & \textbf{要做的事} \\
\midrule
赛后到冬训 & 复盘当年的问题，定下一年的技术路线；把代码整理到能交接 \\
\addlinespace
年初 & 关注报名与资格审查通知；准备 TDP \\
\addlinespace
春季 & 提交材料；按新规则改代码；机器人体检 \\
\addlinespace
赛前两月 & 整机联调；按比赛场景跑全流程；网络和裁判盒演练 \\
\addlinespace
比赛 & 现场调试；每天赛后复盘，当天改的东西当天记录 \\
\bottomrule
\end{tabularx}
\caption{一个赛季中常见的准备阶段与对应任务。}
\end{table}

\subsection{队内信息维护流程}

订阅本身不能保证重要通知得到处理。队内应明确维护责任，并安排固定的检查频率：

\begin{itemize}[leftmargin=*]
    \item 一人盯邮件列表，每周扫一次，规则相关的转进队群。
    \item 一人盯 Discord 的规则与赛事频道，其余频道不用看。
    \item 规则仓库每两周对一次 \texttt{git log}，有改动就记进变更表。
\end{itemize}

\noindent 每个赛季结束后，应保存当届规则版本及对应的 commit 哈希或 release 标签。这样在需要追溯条款变化时，可以定位到具体版本。\advice

\newpage
% ==========================================
% 第 6 章：安全守则
% ==========================================
\section{实验室安全守则}

人形机器人包含高扭矩关节以及高电压、大电流的供电系统。即使机器人暂时没有移动，通电关节也可能突然施力，因此上电、测试、搬运和故障处置都需要遵循固定流程。具体风险应结合平台参数和实验室安全规程评估。

以下流程用于降低误操作和设备失控的风险。任何一项无法确认时，都应暂停上电并联系实验室负责人。

\subsection{上电之前：六项检查}

\begin{enumerate}[leftmargin=*]
    \item \textbf{吊机固定好了吗？} 保护绳受力，留 5 到 10 厘米余量。
    \item \textbf{关节处于规定位置吗？} 按平台说明确认关节姿态，避免在扭曲或机械死点状态下上电。
    \item \textbf{周围清空了吗？} 一米内没有人，地面没有线缆和工具。
    \item \textbf{急停在手边吗？} 按得到，够得着，知道按下去会怎样。
    \item \textbf{遥控器在谁手上？} 指定一个人，拇指放在阻尼键上。
    \item \textbf{三个人都在吗？} 缺一个就不上电。
\end{enumerate}

\noindent 任一项未确认时，不得上电。

\subsection{保护吊架的使用}

保护吊架用于限制机器人失衡后的下坠，不应持续牵引机器人。吊绳过紧会改变测试时的受力条件，影响步态表现。绳长、挂点和允许的松弛量应按平台说明及实验室规程设置。

\begin{itemize}[leftmargin=*]
    \item 绳子调到机器人站直时略松，留 5 到 10 厘米下沉量。
    \item 挂点在躯干重心正上方，不要在头部或者单侧肩膀。
    \item 走路测试时，吊架操作员应跟随机器人移动，并避免主动牵引机器人；吊架只承担规程规定的保护作用。
    \item 测试期间应始终保持保护绳连接。失控跌落可能损坏关节、齿轮箱或结构件。
\end{itemize}

\subsection{三人测试分工}

\begin{table}[H]
\centering
\small
\begin{tabularx}{\textwidth}{p{2.2cm} p{1.4cm} X}
\toprule
\textbf{岗位} & \textbf{代号} & \textbf{职责与站位} \\
\midrule
推吊车 & A & 双手扶吊车，视线盯着机器人重心，跟着移动。随时准备往上提。 \\
\addlinespace
遥控手 & B & 握遥控器，拇指常在阻尼键上。只盯机器人，不聊天。 \\
\addlinespace
监护 & C & 站在急停和电源旁，管人、管线、管记录。出事先拍急停。 \\
\bottomrule
\end{tabularx}
\caption{机器人测试期间三名成员的职责分工。}
\end{table}

\noindent 测试时，吊架操作员和遥控手都可能无法立即处理电源或人员隔离，因此需要独立的监护人负责急停、周边安全和异常记录。

\subsection{阻尼与急停的使用边界}

\begin{table}[H]
\centering
\small
\begin{tabularx}{\textwidth}{p{2.6cm} X X}
\toprule
& \textbf{进阻尼} & \textbf{拍急停} \\
\midrule
做了什么 & 依平台实现切换到阻尼状态，具体电气行为以厂家说明为准 & 通过硬件回路切断动力输出；机器人仍可能因重力或惯性继续运动 \\
\addlinespace
作用 & 在驱动与通信正常时，可能减少主动输出并降低部分冲击 & 可在主控制软件失效时切断动力输出，具体效果取决于硬件设计 \\
\addlinespace
限制 & 依赖控制链路和驱动器仍能响应，不适用于所有故障 & 切断动力可能使机器人失去主动支撑；应按设备设计评估后果 \\
\addlinespace
使用条件 & 按厂家说明和已验证的本队流程决定 & 存在人员危险、烟雾或火花，或其他需要立即切断动力的情况 \\
\bottomrule
\end{tabularx}
\caption{阻尼与急停的作用不同，具体处置顺序应遵循平台说明和实验室批准的流程。}
\end{table}

阻尼模式并不适用于所有故障，也不能替代硬件急停。应根据故障类型、人员位置和设备说明执行经过验证的停机流程；若存在人员受伤风险、出现烟雾或火花，或阻尼指令无响应，应立即按实验室应急流程使用硬件急停。实际处置顺序必须由 T4 厂家手册与实验室安全负责人确认。

\subsection{三类常见危险状态}

\begin{table}[H]
\centering
\small
\begin{tabularx}{\textwidth}{p{2.6cm} p{4.6cm} X}
\toprule
\textbf{类型} & \textbf{什么样子} & \textbf{为什么危险} \\
\midrule
自激振荡 & 全身高频抖动，幅度越来越大 & 控制环路相位错了，越纠越偏，几秒内能把关节甩到限位 \\
\addlinespace
单腿支撑失控 & 一条腿突然伸直或猛收 & 重心瞬间偏移，机器人朝一侧砸，几十公斤砸下来 \\
\addlinespace
线束缠绕 & 转身时线被拉住，机器人僵住 & 电机还在出力，线被扯断，或者把机器人拽倒 \\
\bottomrule
\end{tabularx}
\caption{测试过程中需要识别的三类危险状态。}
\end{table}

\noindent 出现上述现象时，应立即按实验室规定执行停机流程。不要在机器人仍可能驱动关节时靠近检查，也不要把查看日志或关闭软件作为首要处置动作。

\subsection{机器人跌倒后的处理}

\begin{enumerate}[leftmargin=*]
    \item 确认周围人员安全后，按平台规定的应急流程停机并隔离动力源。不得在关节仍可能主动施力时扶起机器人。
    \item 拍照。记录姿态和现场，事后分析靠这个。
    \item 检查三个地方：关节有没有异响、结构件有没有裂纹、线束有没有被扯。
    \item 只有在动力源已隔离、设备负责人确认可以操作后，才检查关节是否卡滞；不要强行转动可能受损的关节。
    \item 重新上电前，应先完成结构、电气和关节检查。首次测试应在保护措施到位的条件下，从低风险动作开始。
\end{enumerate}

\subsection{开机与关机顺序}

\begin{itemize}[leftmargin=*]
    \item \textbf{开机：} 低压先通（工控机、传感器），确认自检通过，喊一声「注意上高压」，再合动力电。
    \item \textbf{关机：} 按平台规定的停机顺序让机器人进入安全状态，确认机械支撑可靠后，再断开动力电并关闭工控机。
    \item \textbf{严禁}在关节扭曲、踩到异物、或者人还扶着机器人的时候合闸。
\end{itemize}

\vspace{0.4cm}
\begin{dangerbox}[贴在工位上的那张卡]
\textbf{上电前：} 吊机就位 \(\cdot\) 关节归零 \(\cdot\) 清场 \(\cdot\) 急停在手 \(\cdot\) 遥控有人 \(\cdot\) 三人到齐\\
\textbf{出事时：} 按实验室批准的应急流程停机；存在人员危险、烟雾或火花时立即使用硬件急停；隔离动力源后再接触机器\\
\textbf{原则：} 先确保人员安全，再隔离动力源并记录设备状态。
\end{dangerbox}

\newpage
% ==========================================
% 第 7 章：第一次参赛
% ==========================================
\section{第一次参赛：从报名到上场}

参赛工作除了研发，还包括资格材料、运输、现场检录和比赛流程。任何一项准备不足都可能影响参赛，即使机器人本身已经能够运行。

\subsection{报名要准备什么}

\begin{table}[H]
\centering
\small
\begin{tabularx}{\textwidth}{p{2.8cm} X p{2.6cm}}
\toprule
\textbf{材料} & \textbf{写什么，容易漏什么} & \textbf{提前多久} \\
\midrule
TDP & 硬件方案、软件架构、感知与决策、实验结果。容易漏的是「和去年比改了什么」 & 2 个月 \\
\addlinespace
资格视频 & 展示自主寻球、带球、射门和起身；按当届资格说明拍摄并确保动作可核验 & 1 个月 \\
\addlinespace
机器人参数表 & 尺寸、重量、自由度、传感器清单 & 1 个月 \\
\addlinespace
队伍信息 & 成员、指导老师、所属机构 & 随表提交 \\
\bottomrule
\end{tabularx}
\caption{报名与资格审查阶段通常需要准备的材料。具体要求和截止日期以当届通知为准。}
\end{table}

资格视频应清楚展示机器人的自主行为。拍摄时可以保留机器人、场地和操作人员的全景，确保评审能够判断动作是否由机器人自主完成。具体拍摄要求以当届资格说明为准。

\subsection{出发前：打包清单}

\begin{table}[H]
\centering
\small
\begin{tabularx}{\textwidth}{p{2.6cm} X}
\toprule
\textbf{类别} & \textbf{清单} \\
\midrule
机器人 & 整机 \(\times\)2（能带就带备机）、备用关节、备用电池、充电器 \\
\addlinespace
工具 & 内六角全套、螺丝、扎带、热风枪、万用表、焊台 \\
\addlinespace
线材 & 网线、交换机、电源延长线、转换插头 \\
\addlinespace
计算 & 调试笔记本、备用硬盘、装好环境的 U 盘 \\
\addlinespace
耗材 & 备用脚底、胶带、防滑垫、清洁剂 \\
\addlinespace
文件 & 规则打印件、TDP 打印件、保险与签证材料 \\
\bottomrule
\end{tabularx}
\caption{出发前需要核对的设备、工具、耗材和文件。}
\end{table}

\subsection{抵达赛场后的检查}

抵达场地后，建议先完成以下检查，再进行整机测试：

\begin{enumerate}[leftmargin=*]
    \item \textbf{确认网络。} 了解赛场网段、局域网限制和网络隔离要求，并确认裁判盒、机器人与调试机之间的通信方式符合现场规定。
    \item \textbf{检查视觉参数。} 在实际场地检查曝光、白平衡和检测效果。实验室与赛场的光照条件不同，是否需要重新标定应由现场测试决定。
    \item \textbf{检查比赛状态切换。} 使用裁判盒依次测试 INITIAL、READY、SET 和 PLAYING 等状态，确认机器人能够正确接收并响应状态变化。
    \item \textbf{复核机器人尺寸。} 使用现场允许的量具确认尺寸、重量与报名材料一致，并留出处理测量差异的时间。
\end{enumerate}

\subsection{比赛日流程}

\begin{table}[H]
\centering
\small
\begin{tabularx}{\textwidth}{p{2.4cm} X}
\toprule
\textbf{时间点} & \textbf{做什么} \\
\midrule
赛前 2 小时 & 充满电、检查结构、跑一次完整流程 \\
\addlinespace
赛前 1 小时 & 换好电池、确认队伍编号、和裁判对一遍规则 \\
\addlinespace
赛前 15 分钟 & 机器人进场、摆位、最后一次状态检查 \\
\addlinespace
比赛中 & 记录异常现象及发生时间；除非确有必要且留有验证时间，否则避免修改代码 \\
\addlinespace
赛后 30 分钟 & 复盘、改一处、重新验证，当天记录 \\
\bottomrule
\end{tabularx}
\caption{比赛日各时间点的准备与检查事项。}
\end{table}

\subsection{现场调试与变更记录}

\begin{enumerate}[leftmargin=*]
    \item \textbf{一次只改一处。} 同时改三个参数，好了不知道哪个起作用，坏了也不知道哪个要退。
    \item \textbf{改了就记。} 记下时间、改了什么、现象怎么变。赛后复盘全靠这本流水账。
    \item \textbf{能回滚。} 改之前提交一次，或者至少备份文件。现场没有时间重写。
\end{enumerate}

\noindent 可以将赛前两小时设为本队的变更冻结点：此后原则上只调整经过验证的参数，不再改动软件架构。若必须修复严重故障，应先评估剩余验证时间和回滚方案。\advice

\subsection{赛后复盘模板}

建议在比赛当天完成记录，避免依赖几天后可能不准确的回忆：

\begin{itemize}[leftmargin=*]
    \item 这场输在哪：丢球、误判、走位、还是没电？
    \item 哪个现象是第一次出现？之前有没有预兆？
    \item 如果只改一处，改哪里？
    \item 下一场之前必须验证什么？
\end{itemize}

\noindent 复盘记录应尽量对应到可观察的现象、可能原因和下一步验证任务。这样后续修改才有依据，也便于下一届队员了解问题是如何解决的。

\newpage
% ==========================================
% 附录
% ==========================================
\appendix
\section{术语表}

\begin{longtable}{p{2.6cm} p{4.2cm} p{7.4cm}}
\toprule
\textbf{中文} & \textbf{英文 / 缩写} & \textbf{解释} \\
\midrule
\endfirsthead
\toprule
\textbf{中文} & \textbf{英文 / 缩写} & \textbf{解释} \\
\midrule
\endhead
\bottomrule
\endfoot
人形足球联赛 & HSL & RoboCup 下属的人形机器人足球联赛，本书的主角 \\
\addlinespace
标准平台组 & SPL & 统一使用 NAO 的历史组别，已并入人形体系 \\
\addlinespace
中型组 & MSL & 轮式中型机器人组别，已停办 \\
\addlinespace
技术文档 & TDP & 各队每年提交的技术说明，含方案与结果 \\
\addlinespace
裁判盒 & GC & 广播比赛状态的官方程序，见 1.6 节 \\
\addlinespace
数据分发服务 & DDS & ROS 2 底层的通信标准，无中心节点 \\
\addlinespace
服务质量 & QoS & 规定消息丢了要不要重传、队列多深 \\
\addlinespace
坐标变换 & TF2 & 维护各坐标系之间变换的 ROS 2 组件 \\
\addlinespace
坐标变换树 & TF tree & 从地图到相机的一串变换关系 \\
\addlinespace
蒙特卡洛定位 & MCL & 用粒子表示位姿分布的定位方法，见 4.7 节 \\
\addlinespace
粒子 & particle & 一个候选位姿，MCL 的基本单位 \\
\addlinespace
重采样 & resampling & 按权重复制和淘汰粒子的一步 \\
\addlinespace
里程计 & odometry & 靠关节和触地推算的相对位移 \\
\addlinespace
正运动学 & FK & 已知关节角，算脚在哪 \\
\addlinespace
逆运动学 & IK & 已知脚要去哪，算关节角 \\
\addlinespace
惯性测量单元 & IMU & 测角速度和加速度的传感器 \\
\addlinespace
零力矩点 & ZMP & 判定机器人会不会倒的一个点 \\
\addlinespace
压力中心 & CoP & 足底受力的等效作用点 \\
\addlinespace
有限状态机 & FSM & 状态加跳转条件的控制结构 \\
\addlinespace
行为树 & BT & 用组合节点组织行为的控制结构，见 4.6 节 \\
\addlinespace
仿真到真机 & sim2real & 把仿真里训好的策略搬到真机 \\
\addlinespace
域随机化 & domain randomization & 训练时随机改参数，逼策略适应差异 \\
\addlinespace
机器人描述文件 & URDF & 描述连杆、关节、外观的模型文件 \\
\addlinespace
推理引擎 & TensorRT / .engine & 把模型编译成显卡能快速跑的格式 \\
\addlinespace
漏检 & false negative & 场景中存在目标，但检测系统未能识别该目标 \\
\addlinespace
误检 & false positive & 场景中不存在目标，但检测系统输出了目标结果 \\
\end{longtable}

\section{新成员 90 天计划}

这份计划按每周约 8 至 10 小时安排，目标是在三个月内让新成员能够独立修改并验证一个小模块。实际进度应根据课程负担和队内任务调整。

\begin{table}[H]
\centering
\small
\begin{tabularx}{\textwidth}{p{1.8cm} p{3.2cm} X}
\toprule
\textbf{周次} & \textbf{做什么} & \textbf{验收标准} \\
\midrule
第 1--2 周 & 读第 1、6 章；看一次完整的真机测试 & 能说清比赛状态和三种安全操作 \\
\addlinespace
第 3--4 周 & Git 与 Linux；克隆仓库，本地编译一遍 & 能提一个 PR，能改一个参数 \\
\addlinespace
第 5--6 周 & 第 4.2、4.3 节；跑通 ROS 2 示例 & 能画出队里的话题拓扑 \\
\addlinespace
第 7--8 周 & 跟着做一次真机联调；学录包和回放 & 能用 rosbag2 复现一个现象 \\
\addlinespace
第 9--10 周 & 进组：感知读 4.4，策略读 4.5，运控读 3.2 & 能在自己的模块里改一个阈值并解释结果 \\
\addlinespace
第 11--12 周 & 独立接一个小任务（比如换个检测类别） & 从改代码到上机验证，全程自己走一遍 \\
\bottomrule
\end{tabularx}
\caption{按每周约 8 至 10 小时设计的新成员学习计划。}
\end{table}

\section{故障排查速查}

\begin{longtable}{p{3.4cm} p{5.2cm} p{5.6cm}}
\toprule
\textbf{现象} & \textbf{先查什么} & \textbf{常见原因} \\
\midrule
\endfirsthead
\toprule
\textbf{现象} & \textbf{先查什么} & \textbf{常见原因} \\
\midrule
\endhead
\bottomrule
\endfoot
机器人完全不动 & 裁判盒状态、急停、阻尼开关 & 状态还在 INITIAL；急停没复位 \\
\addlinespace
收不到裁判盒 & 队号、球员号、网段、防火墙 & 编号不一致；广播被拦 \\
\addlinespace
定位乱跳 & 特征数量、粒子数、重采样阈值 & 镜像解；特征太少还强行重采样 \\
\addlinespace
检测不到球 & 曝光、阈值、模型输入尺寸 & 过曝；训练集没有这个场地 \\
\addlinespace
画面延迟大 & QoS 配置、队列深度、CPU 占用 & 可靠传输加深度队列 \\
\addlinespace
步态发散 & IMU 零偏、关节零位、吊绳是否绷紧 & 零偏没减；吊绳拉着走 \\
\addlinespace
关节异响 & 立即断电，手动转关节 & 齿轮损伤；线束干涉 \\
\addlinespace
走直线偏 & 左右腿零位差、地面摩擦 & 零位没有对正；一侧打滑 \\
\addlinespace
踢球方向偏 & 髋旋转零位、球的位置估计 & 零位偏差；视觉坐标没转对 \\
\addlinespace
编译失败 & 依赖版本、工作空间是否 source & 环境没加载；分支不对 \\
\addlinespace
上电无反应 & 电池电压、急停回路、保险 & 电池过放；急停没复位 \\
\addlinespace
通信间歇丢包 & 网线、交换机、无线干扰 & 网线松动；多人抢同一频段 \\
\end{longtable}

\section{待补事实清单}

本节汇总正文中标记为 \todo{待补} 的信息。补齐并核对后，可删除本节。

\begin{table}[H]
\centering
\small
\begin{tabularx}{\textwidth}{p{3.6cm} p{6.4cm} X}
\toprule
\textbf{条目} & \textbf{去哪拿} & \textbf{谁负责} \\
\midrule
T4 整机尺寸与重量 & 厂家手册 / 自己量 & \\
\addlinespace
关节限位表 & SDK 文档 + 实测 & \\
\addlinespace
电池续航与更换策略 & 实测记录 & \\
\addlinespace
队史与历年名次 & 队内档案 & \\
\addlinespace
网站平台功能清单 & 平台负责人 & \\
\addlinespace
规则条款编号 & Rules.tex 对应小节 & \\
\addlinespace
各队近年战绩 & 官方 qualified-teams & \\
\bottomrule
\end{tabularx}
\caption{需要由队内负责人补齐并定期核验的信息。}
\end{table}

\clearpage
\section{队内黑话对照}

下表列出队内常用说法及其含义，便于新成员阅读文档和参与讨论。

\begin{table}[H]
\centering
\small
\begin{tabularx}{\textwidth}{p{2.8cm} X}
\toprule
\textbf{说法} & \textbf{意思} \\
\midrule
上机 / 上真机 & 在实体机器人上运行；仿真环境里不算 \\
\addlinespace
跑包 & 用 rosbag2 回放录好的数据 \\
\addlinespace
挂机 & 机器人吊在吊机上，准备测试 \\
\addlinespace
进阻尼 & 切到阻尼模式，关节失去主动力（见第 6 章） \\
\addlinespace
翻车 & 机器人摔倒，或者整套流程失效 \\
\addlinespace
对零位 & 把关节摆到标定零位 \\
\addlinespace
标场地 & 在真实场地上重新调相机参数 \\
\addlinespace
踢空 & 检测到球但位置算错，出脚落空 \\
\addlinespace
追不上 & 决策给了目标速度，底层的实际速度跟不上 \\
\addlinespace
掉包 & 话题或裁判盒数据丢了 \\
\addlinespace
烧板子 & 驱动板损坏，通常伴随焦味 \\
\addlinespace
过检录 & 尺寸、重量、传感器全部合规 \\
\bottomrule
\end{tabularx}
\caption{队内常见术语与口头表达。}
\end{table}

\section{延伸阅读}

\begin{table}[H]
\centering
\small
\begin{tabularx}{\textwidth}{p{3.4cm} X}
\toprule
\textbf{资源} & \textbf{读它解决什么问题} \\
\midrule
Google 开发者文档风格指南 & 关于清晰表达、准确用词和避免过度断言的写作建议 \\
\addlinespace
ROS 2 官方文档 & 概念定义清楚，优先读 Concepts 部分 \\
\addlinespace
B-Human Code Release 文档 & 一套完整机器人系统的架构说明，见第 2 章 \\
\addlinespace
HSL-Rules 仓库 & 规则原文与修订历史。有歧义时以它为准 \\
\addlinespace
BehaviorTree.CPP 文档 & 行为树的工程实现，含可视化工具 \\
\addlinespace
Ultralytics YOLO 文档 & 训练、导出、部署的完整流程 \\
\addlinespace
Linux Journey & 命令行入门，做到「文件、进程、权限」三章即可 \\
\bottomrule
\end{tabularx}
\caption{可用于继续学习的官方文档、项目资料和工具说明。}
\end{table}
\vspace{0.6cm}
\begin{center}
\rule{0.6\textwidth}{0.5pt}\\[0.3cm]
{\small \textsf{Robot 工作坊 \(\cdot\) BigHeroX 战队研发委员会 编制}}\\[0.1cm]
{\footnotesize \textsf{本书对应 2026 赛季 \(\cdot\) 规则以 HSL-Rules 仓库为准}}\\[0.1cm]
{\footnotesize \textsf{All Rights Reserved \(\copyright\) 2026}}
\end{center}

\end{document}
